\section{Introducci\'on}

\subsection{Motivaci\'on}

El siguiente paquete ofrece una colección de macros diseñadas para representar tableros y posiciones de Go, implementadas íntegramente en \pgfname. Este proyecto, aunque aún en desarrollo, incorpora funcionalidades innovadoras que buscan simplificar y agilizar la creación de diagramas relacionados con el juego de Go:
\begin{itemize}
	\item Gestión de perfiles para tableros y jugadores.
	\item Carga de múltiples coordenadas simultáneamente mediante una sintaxis que algunos podrían encontrar familiar.
\end{itemize}

El desarrollo de \pgfgoname\ surge de la frustración por no encontrar paquetes que satisficieran mis necesidades. Los paquetes que exploré permiten representar tableros en blanco y negro, adecuados para impresión, pero poco prácticos al combinarlos con tecnolog\'ias como |beamer| o en documentos donde se desea una presentación a color.

Inicialmente, \pgfgoname\ se limitaría a estas tareas; sin embargo, ahora lo concibo más como una interfaz para gestionar estilos y estados, evitando el uso excesivo del formato directo. Esto se reflejará principalmente en la estructura del árbol de \emph{keys}.

En cuanto al nombre, aunque las macros están implementadas en \pgfname\ y L3, el tablero se construye sobre |\tikzpicture| para que los usuarios familiarizados con este entorno puedan aprovechar sus macros.


\subsection{Novedades}


\begin{itemize}
	\item \pgfgoversion:
		\begin{itemize}
			\item El cálculo (coordenadas, numeración de jugadas, memoria, etc.) se ha separado del dibujo. Lo hace un \emph{kernel}, que se elige con la opción |kernel| del paquete (pág.~\pageref{subsec:kernel}). Además del kernel escrito en L3, se ha añadido un kernel en Lua, que es el que se usa por defecto con Lua\LaTeX.
			\item Con el kernel en Lua se han añadido las capturas (|captures|, pág.~\pageref{subsec:capturas}) y la lectura de partidas en SGF (|\sgf|, pág.~\pageref{subsec:sgf}). Con el paquete |animate|, |\sgfanimation| permite recorrer una partida con botones.
			\item Se han añadido las notas «N en M» para las jugadas numeradas en puntos ya jugados (|notes| y |\pgfgomovenotes|, pág.~\pageref{subsec:notas}).
			\item SGF: variantes (|variation|), marcas del archivo (|marks|), los datos de la partida (|\sgfinfo|), la numeración de |MN| y |restart|, y el tamaño del tablero tomado del archivo (|sgf size|).
			\item SGF: colecciones de partidas (|game|), la vista (|VW|, |sgf view|), las figuras (|FG|) en |\sgfdiagrams|, las variantes como letras (|variations|), |PL| y los archivos en Latin-1.
			\item SGF: el nombre de los nodos y las anotaciones (|TE|, |BM|, |GB|, etc.) se escriben con los comentarios. |\sgfcomment| escribe el comentario de una jugada, y |\sgfdiagrams| puede escribir los comentarios debajo de cada diagrama (|comments|).
			\item Se han añadido flechas y líneas (|\arrow| y |\line|, pág.~\pageref{subsec:flechas}) y la marca |dim|. |\sgf| dibuja además |AR|, |LN|, |SL| y |DD|.
			\item El tablero dibuja los puntos de estrella (|hoshi|, pág.~\pageref{subsec:hoshi}).
			\item Se han añadido |\territory|, que marca el territorio (con un polígono, o solo con |auto| y el kernel en Lua), y |\pgfgoscore|, que cuenta los puntos.
			\item |\sgfanimation| arma la animación por capas (|layers|): cada piedra se dibuja una sola vez, y el PDF de una partida larga ocupa mucho menos.
			\item Con |captured marks = remove|, las marcas sobre una piedra capturada se van con ella.
			\item Se han añadido |\pgfgoprisoners|, que cuenta los prisioneros de cada jugador, y |\sgfdiagrams|, que dibuja una partida como serie de diagramas, como en los libros.
			\item Se ha añadido |\sgfproblem|, que dibuja un problema (tsumego): el enunciado y un diagrama por cada variante de la solución, con sus comentarios.
			\item SGF: la numeración que pide |PM| (|numbers| y |restart = auto|), y la vista (|VW|) cuando cambia durante la partida (|sgf view move|; |\sgfdiagrams| y |\sgfanimation| la usan solos). |\sgfdiagrams| escribe el resultado y el puntaje debajo del último diagrama (|result|).
			\item Las coordenadas del borde pueden ser letras, números o kanji en cada eje, contados desde cualquier lado (|columns|, |rows|, |style = japanese|, etc.).
		\end{itemize}
	\item 0.2025.10.13: Corrección de errores. Cada |goban| empieza en la jugada 1 con negras, |\mark| y |\marks| solo existen dentro de |goban|, y lo que cae fuera de un tablero parcial ya no se dibuja. La lista completa está en el archivo |CHANGELOG.md|.
	\item 0.2025.03.10:
		\begin{itemize}
			\item Se han a\~nadido las \emph{keys} |inner| y |outer| para el \emph{rounded corner} del tablero.
			\item Se han a\~nadido tres nuevos operadores de coordenadas: |*| para todas las intersecciones, |L| para toda la columna de coordenadas L y |N| para toda la fila de coordenadas N.
		\end{itemize}
	\item 0.2025.03.07: Primer lanzamiento.
\end{itemize}